\chapter{Authoring Guide}
\label{ch:authoring-guide}

This chapter teaches you to write a \texttt{gert.ops} runbook from scratch. It walks through
role declaration, approval gates, change-request wrapping, and evidence capture, ending with
a complete, realistic deployment runbook you can use as a template.

% ─────────────────────────────────────────────────────────────────────────────
\section{Starting a New Runbook}
\label{sec:authoring-start}
% ─────────────────────────────────────────────────────────────────────────────

Every \texttt{gert.ops} runbook begins with three declarations: the schema version, the Kit,
and the metadata block.

\begin{minted}{yaml}
apiVersion: runbook/v2
kit: ops/v1

meta:
  name: deploy-api-service
  kind: change
  description: >
    Deploy a new version of the API service to production.
    Requires change-board approval before execution.
  inputs:
    service_name:
      from: prompt
      description: "Service identifier (e.g., api-gateway)"
    image_tag:
      from: prompt
      description: "Docker image tag to deploy (e.g., v2.4.1)"
    change_ticket:
      from: prompt
      description: "Change request ticket number (e.g., CHG0042817)"
    previous_image_tag:
      from: prompt
      description: "Previous image tag, used for rollback"
\end{minted}

\paragraph{Choosing a \texttt{kind}.}
Use \texttt{change} for planned deployments and infrastructure modifications, \texttt{incident}
for incident-response runbooks, and \texttt{operational} for routine operational procedures
(e.g., on-call runbooks, diagnostic playbooks).

% ─────────────────────────────────────────────────────────────────────────────
\section{Declaring Roles}
\label{sec:authoring-roles}
% ─────────────────────────────────────────────────────────────────────────────

Declare all participants in \texttt{meta.roles}. Be explicit: role declarations are enforced
at runtime. A step with \texttt{requires\_role: dri} will be blocked if the executing user
does not hold the \texttt{dri} role.

\begin{minted}{yaml}
meta:
  roles:
    dri: "@alice"
    approver:
      - "@bob"
      - "@change-board"
    change-manager: "@change-board"
    observer:
      - "@audit-log-team"
  sla:
    approval-timeout: 24h
    escalation-target: "@sre-lead"
\end{minted}

\paragraph{Rule of thumb.} Start with the DRI (who is accountable?), then the approver
(who must sign off?), and the change-manager (who authorized the change window?). Add
observers for audit requirements. If you are writing an incident runbook, also declare
\texttt{responder} and \texttt{incident-commander}.

% ─────────────────────────────────────────────────────────────────────────────
\section{Governance Configuration}
\label{sec:authoring-governance}
% ─────────────────────────────────────────────────────────────────────────────

Declare the command allowlist and output redaction rules in \texttt{meta.governance}.
For a deployment runbook, this typically covers the tools you use:

\begin{minted}{yaml}
meta:
  governance:
    allowed_commands: [kubectl, docker, curl, jq, helm]
    denied_commands: [rm, dd, shutdown]
    deny_env_vars:
      - "*_SECRET*"
      - "*_TOKEN"
      - "AWS_SECRET_ACCESS_KEY"
    redact:
      - pattern: "(--password=)[^ ]+"
        replace: "$1[REDACTED]"
      - pattern: "Bearer [A-Za-z0-9._-]+"
        replace: "Bearer [REDACTED]"
\end{minted}

Every \texttt{ops.cli} step in the runbook inherits these rules. A step may additionally
declare its own \texttt{allowed\_commands} list, which is intersected with the runbook-level
list.

% ─────────────────────────────────────────────────────────────────────────────
\section{Using Approval Gates}
\label{sec:authoring-approval}
% ─────────────────────────────────────────────────────────────────────────────

\paragraph{When to use an approval gate.}
Use \texttt{ops.approval} whenever execution must pause for a human decision:
\begin{itemize}
  \item Before making a change to a production environment.
  \item After evidence has been gathered and a reviewer needs to sign off before continuing.
  \item At any point where policy requires documented authorization.
\end{itemize}

\paragraph{When \emph{not} to use an approval gate.}
Do not use approval gates for informational pauses (use \texttt{ops.manual} without evidence
instead). Do not chain multiple gates on the same step sequence unless they represent distinct
approval authorities (e.g., a technical lead \emph{and} a change manager are different roles).

\begin{minted}{yaml}
flow:
  - step:
      id: request_approval
      type: ops.approval
      title: "Change board approval required"
      description: |
        Deployment: {{ .service_name }} → {{ .image_tag }}
        Change ticket: {{ .change_ticket }}
        Planned window: {{ .change_window }}
      approvers:
        - "@change-board"
      requires_any: true
      timeout: 24h
      on_timeout: escalate
      escalation_target: "@vp-engineering"
      evidence:
        - type: ops.evidence.attestation
          label: "Change board sign-off"
          required: true
          prompt: "I approve this deployment for the scheduled change window."
\end{minted}

\paragraph{Timeout behavior.} When the timeout fires:
\begin{itemize}
  \item With \texttt{on\_timeout: escalate}: The run is suspended. The
    \texttt{escalation\_target} receives a notification. The run can be resumed once approval
    is granted or the gate is bypassed by the DRI with documented justification.
  \item With \texttt{on\_timeout: abort}: The run is cancelled. A trace event is written.
    No further steps execute.
\end{itemize}

% ─────────────────────────────────────────────────────────────────────────────
\section{Using the Change-Request Wrapper}
\label{sec:authoring-change-request}
% ─────────────────────────────────────────────────────────────────────────────

Wrap your deployment or rollback execution steps in \texttt{ops.change-request} to activate
the change management lifecycle. This provides automatic pre-execution approval from the
change-manager, rollback injection on failure, and change record closure.

\begin{minted}{yaml}
  - step:
      id: execute_change
      type: ops.change-request
      title: "Deploy {{ .service_name }} {{ .image_tag }}"
      change-id: "{{ .change_ticket }}"
      risk: medium
      rollback:
        runbook: ./rollback-api-service.runbook.yaml
        inputs:
          service_name: "{{ .service_name }}"
          target_image_tag: "{{ .previous_image_tag }}"
      steps:
        - step:
            id: apply_deployment
            type: ops.cli
            title: "Apply deployment manifest"
            command: kubectl
            args: [set, image, "deployment/{{ .service_name }}",
                   "app={{ .image_tag }}", "-n", production]
            evidence:
              auto: true
              label: "kubectl set image output"
        - step:
            id: wait_rollout
            type: ops.cli
            title: "Wait for rollout to complete"
            command: kubectl
            args: [rollout, status, "deployment/{{ .service_name }}",
                   "-n", production, "--timeout=5m"]
            evidence:
              auto: true
              label: "rollout status output"
\end{minted}

\textbf{Key point:} The \texttt{ops.change-request} wrapper does not execute the change
automatically. The change-manager must first approve the pre-execution gate (populated from
\texttt{meta.roles.change-manager}). Only after that approval do the \texttt{steps} execute.

% ─────────────────────────────────────────────────────────────────────────────
\section{Capturing Evidence}
\label{sec:authoring-evidence}
% ─────────────────────────────────────────────────────────────────────────────

Choose the evidence type that best matches the proof you need:

\begin{description}
  \item[\texttt{ops.evidence.command-output}] Use when the proof is the output of a command
    (e.g., \texttt{kubectl rollout status}, \texttt{curl /healthz}). Enable
    \texttt{evidence.auto: true} on \texttt{ops.cli} steps to capture automatically.
  \item[\texttt{ops.evidence.screenshot}] Use when the proof is a visual dashboard state
    (e.g., Grafana, Datadog) that cannot be captured programmatically. Requires the executor
    to upload a file.
  \item[\texttt{ops.evidence.attestation}] Use for formal sign-offs: DRI acknowledgements,
    approvals, and change-manager authorizations. The executor affirms a specific statement.
\end{description}

\paragraph{Evidence on manual steps.}
\begin{minted}{yaml}
  - step:
      id: health_check
      type: ops.manual
      title: "Confirm service health"
      requires_role: dri
      instructions: |
        1. Open Grafana: https://grafana.internal/d/api-service
        2. Confirm error rate < 0.1% for 5 minutes post-deploy.
        3. Confirm p99 latency < 200ms.
        4. Take a screenshot of the dashboard.
      evidence:
        - type: ops.evidence.screenshot
          label: "Grafana health dashboard"
          required: true
          description: "Dashboard showing error rate and latency, 5 min post-deploy."
        - type: ops.evidence.attestation
          label: "DRI health sign-off"
          required: true
          prompt: "I confirm the service is healthy and the deployment is successful."
\end{minted}

% ─────────────────────────────────────────────────────────────────────────────
\section{Best Practices}
\label{sec:authoring-best-practices}
% ─────────────────────────────────────────────────────────────────────────────

\begin{description}
  \item[One DRI, always.] Every runbook must have a clear DRI. If the runbook can be run
    by anyone, use \texttt{dri: "@oncall-sre"} to assign the on-call engineer.
  \item[Gate before irreversible actions.] Place an \texttt{ops.approval} gate before every
    step sequence that modifies production state and cannot be trivially undone.
  \item[Always specify rollback.] Every \texttt{ops.change-request} wrapper must have a
    \texttt{rollback} specification. Do not skip this even for ``low risk'' changes.
  \item[Capture evidence at the point of action.] Use \texttt{evidence.auto: true} on
    \texttt{ops.cli} steps rather than a separate manual evidence step after the fact. The
    evidence should capture the actual state, not a post-hoc reconstruction.
  \item[Keep instructions human-readable.] The \texttt{instructions} field in
    \texttt{ops.manual} steps will be read under pressure. Use numbered lists and concrete
    URLs. Avoid pronouns and ambiguous references.
  \item[Use \texttt{meta.change-id}.] Always link to the external change ticket. This
    allows compliance teams to cross-reference the gert evidence bundle with the change
    management system.
  \item[Test the rollback runbook.] Run the rollback runbook in a non-production environment
    before scheduling the change window. An untested rollback is not a rollback.
\end{description}

% ─────────────────────────────────────────────────────────────────────────────
\section{Complete Example: Production Deployment}
\label{sec:authoring-complete-example}
% ─────────────────────────────────────────────────────────────────────────────

The following is a complete, production-ready deployment runbook demonstrating all major
\texttt{gert.ops} constructs.

\begin{minted}{yaml}
apiVersion: runbook/v2
kit: ops/v1

meta:
  name: deploy-api-service
  kind: change
  description: >
    Deploy a new version of the API service to production.
    Requires change-board approval before execution begins.
  inputs:
    service_name:
      from: prompt
      description: "Service identifier (e.g., api-gateway)"
    image_tag:
      from: prompt
      description: "Docker image tag to deploy (e.g., v2.4.1)"
    change_ticket:
      from: prompt
      description: "Change ticket number (e.g., CHG0042817)"
    previous_image_tag:
      from: prompt
      description: "Previous image tag for rollback"
    namespace:
      from: prompt
      description: "Kubernetes namespace (e.g., production)"
      default: "production"
  roles:
    dri: "@alice"
    approver:
      - "@change-board"
    change-manager: "@change-board"
    observer:
      - "@audit-team"
  sla:
    approval-timeout: 24h
    escalation-target: "@sre-lead"
  change-id: "{{ .change_ticket }}"
  governance:
    allowed_commands: [kubectl, curl, jq]
    deny_env_vars: ["*_SECRET*", "*_TOKEN"]

flow:
  # ── Gate 1: Change board approval ─────────────────────────────────────────
  - step:
      id: change_board_approval
      type: ops.approval
      title: "Change board approval"
      description: |
        Deployment request:
          Service:  {{ .service_name }}
          Image:    {{ .image_tag }}
          Ticket:   {{ .change_ticket }}
          Namespace: {{ .namespace }}
      approvers: ["@change-board"]
      requires_any: true
      timeout: 24h
      on_timeout: escalate
      escalation_target: "@sre-lead"
      evidence:
        - type: ops.evidence.attestation
          label: "Change board approval"
          required: true
          prompt: >
            I approve deployment of {{ .image_tag }} to {{ .namespace }}
            under change ticket {{ .change_ticket }}.

  # ── Change request wrapper ─────────────────────────────────────────────────
  - step:
      id: execute_deployment
      type: ops.change-request
      title: "Deploy {{ .service_name }} {{ .image_tag }}"
      change-id: "{{ .change_ticket }}"
      risk: medium
      rollback:
        runbook: ./rollback-api-service.runbook.yaml
        inputs:
          service_name: "{{ .service_name }}"
          target_image_tag: "{{ .previous_image_tag }}"
          namespace: "{{ .namespace }}"
      steps:
        - step:
            id: set_image
            type: ops.cli
            title: "Update deployment image"
            command: kubectl
            args:
              - set
              - image
              - "deployment/{{ .service_name }}"
              - "app={{ .image_tag }}"
              - "-n"
              - "{{ .namespace }}"
            evidence:
              auto: true
              label: "kubectl set image"

        - step:
            id: wait_rollout
            type: ops.cli
            title: "Wait for rollout to complete"
            command: kubectl
            args:
              - rollout
              - status
              - "deployment/{{ .service_name }}"
              - "-n"
              - "{{ .namespace }}"
              - "--timeout=5m"
            capture:
              rollout_status: stdout
            evidence:
              auto: true
              label: "rollout status"

  # ── Post-deploy health check ───────────────────────────────────────────────
  - step:
      id: health_check
      type: ops.manual
      title: "Verify deployment health"
      requires_role: dri
      instructions: |
        1. Open Grafana: https://grafana.internal/d/{{ .service_name }}
        2. Confirm error rate < 0.1% over 5 min post-deploy.
        3. Confirm p99 latency < 200ms.
        4. Confirm no new alerts in PagerDuty.
        5. Take a screenshot of the Grafana dashboard.
      evidence:
        - type: ops.evidence.screenshot
          label: "Grafana health dashboard"
          required: true
          description: "Dashboard showing error rate and latency, 5 min post-deploy."
        - type: ops.evidence.attestation
          label: "DRI health sign-off"
          required: true
          prompt: >
            I confirm {{ .service_name }} {{ .image_tag }} is healthy in production.
            Error rate is within SLO and latency is acceptable.
\end{minted}
